\begin{table}[H]
\centering
\caption{Síntesis bibliométrica consolidada del corpus (n=73, 2024--2026).}
\label{tab:bibliometria}
\small
\renewcommand{\arraystretch}{1.35}
\begin{tabular}{|>{\raggedright\arraybackslash}m{0.30\linewidth}|>{\raggedright\arraybackslash}m{0.60\linewidth}|}
\hline
\textbf{Dimensión} & \textbf{Resultado} \\
\hline
Evolución de la producción & 2024: 8; 2025: 28; 2026: 37 \\
\hline
Impacto de citación & 40 de 73 estudios (55\%) registran cero citas. La métrica no es interpretable en una ventana de tres años y se reporta como dato descriptivo. \\
\hline
Fuentes de publicación & 54 revistas distintas tras normalización de denominaciones abreviadas. \\
\hline
Zonas de Bradford & Zona 1: 7 revistas / 26 artículos; Zona 2: 25 revistas / 25 artículos; Zona 3: 22 revistas / 22 artículos. \\
\hline
Ajuste a la Ley de Bradford & \textbf{No se cumple.} Los multiplicadores entre zonas son 3.57 y 0.88, lejos de la constancia que exige la ley. \\
\hline
Dispersión de fuentes & 47 de 54 revistas (87\%) aportan un único artículo; esas fuentes concentran el 64\% del corpus. \\
\hline
Núcleo editorial & \emph{Ieee Access} (9); \emph{Ieee Transactions On Mobile Computing} (6); \emph{Visual Intelligence} (3). \\
\hline
Autores & 371 autores únicos: 355 con un artículo (95.7\%), 13 con 2, 2 con 3, 1 con 4. \\
\hline
Ajuste a la Ley de Lotka & La proporción de autores con un solo artículo (95.7\%) excede ampliamente el 62.7\% teórico. La ventana de tres años impide distinguir transitoriedad real de truncamiento del período de observación. \\
\hline
\end{tabular}
\end{table}